\documentclass{article}

\usepackage[margin=1in]{geometry}
\usepackage{graphicx, amsmath, amsthm, amssymb, enumitem, lipsum, hyperref}

\setlist[enumerate]{leftmargin=1.25cm}

\newcommand{\N}{\mathbb{N}}
\newcommand{\Z}{\mathbb{Z}}
\newcommand{\Q}{\mathbb{Q}}
\newcommand{\R}{\mathbb{R}}
\newcommand{\C}{\mathbb{C}}

\newcommand{\rotvert}{\rotatebox[origin=c]{90}{$\vert$}}

\newenvironment{note}[2][Note]{\begin{trivlist}
\item[\hskip \labelsep {\bfseries #1}\hskip \labelsep {\bfseries #2.}]}{\end{trivlist}}

\newenvironment{fact}[2][Fact]{\begin{trivlist}
\item[\hskip \labelsep {\bfseries #1}\hskip \labelsep {\bfseries #2.}]}{\end{trivlist}}

\newenvironment{definition}[2][Definition]{\begin{trivlist}
\item[\hskip \labelsep {\bfseries #1}\hskip \labelsep {\bfseries #2.}]}{\end{trivlist}}

\newenvironment{proposition}[2][Proposition]{\begin{trivlist}
\item[\hskip \labelsep {\bfseries #1}\hskip \labelsep {\bfseries #2.}]}{\end{trivlist}}

\newenvironment{principle}[2][Principle]{\begin{trivlist}
\item[\hskip \labelsep {\bfseries #1}\hskip \labelsep {\bfseries #2.}]}{\end{trivlist}}

\newenvironment{lemma}[2][Lemma]{\begin{trivlist}
\item[\hskip \labelsep {\bfseries #1}\hskip \labelsep {\bfseries #2.}]}{\end{trivlist}}

\newenvironment{theorem}[2][Theorem]{\begin{trivlist}
\item[\hskip \labelsep {\bfseries #1}\hskip \labelsep {\bfseries #2.}]}{\end{trivlist}}

\newenvironment{corollary}[2][Corollary]{\begin{trivlist}
\item[\hskip \labelsep {\bfseries #1}\hskip \labelsep {\bfseries #2.}]}{\end{trivlist}}

\newenvironment{envsection}[1]{\begin{trivlist}
\item[\hskip \labelsep {\bfseries #1}]}{\end{trivlist}}


\begin{document}
\setlength{\abovedisplayskip}{4pt}
\setlength{\belowdisplayskip}{4pt}


\begin{center}
    \Large{\textbf{Homework 6}}\\
    \normalsize March 19 - 26
\end{center}
\vspace{10pt}

\noindent This week, I'm having you mostly focus on readings so you can get as far as you can so we can get to neural nets soon. I won't have as much time myself this week to read the material along with you, so, I thought for our meeting on Thursday, we'd flip things around a bit and have you explain the reading to me on the board. Nothing formal at all, just be prepared to go over the topics with me!

\subsection*{Reading}

\begin{envsection}{1. S+L Processing:}
    Chapter 3: Sections 3.1 (skip 3.1.3), 3.3 (skip 3.3.1), 3.7. You're welcome to look at other sections too, these are just the most important. I want you to look at 3.7 since it talks about a tiny bit of information theory.
\end{envsection}

\begin{envsection}{2. S+L Processing:}
    Chapter 4: Sections 4.1--4.3, 4.5--4.8. You can just skim 4.1, skip 4.3.2, and skip the intro in 4.6 and go straight to 4.6.1. If anything is too familiar from your regression class you're welcome to skim parts.
\end{envsection}

\subsection*{Problems}

\begin{envsection}{Problem 1.}
    Use induction to prove the \textit{chain rule of probability}: For $n$ random variables $X_1,\dots,X_n$,
    \[
    P(X_1,\dots,X_n) = \prod_{k=1}^{n} P(X_k \mid X_{1:k-1}).
    \]
    (Here, $X_{1:k-1}$ refers to the joint distribution $X_1,\dots,X_{k-1}$).
\end{envsection}

\subsection*{Coding}

\begin{envsection}{1.}
    Go back to HW 4 and do some TensorTonic problems I linked.
\end{envsection}

\end{document}